\chapter{Notation and Status of Claims}

This appendix freezes the notation used throughout the book so that one symbol does not quietly change ontological role from chapter to chapter. $\mathscr R$ denotes the underlying relational state space used as the framework's broad mathematical carrier, while $r$ denotes one relational state. $X$ is a raw representation space, $G_{\mathrm{gauge}}$ the group of representational redundancies, $Q=X/G_{\mathrm{gauge}}$ the gauge-reduced state space, and $G_{\mathrm{perm}}$ a problem-specific family of transformations treated as irrelevant to the selected form identity. $\bar Q$ denotes the further reduced space. $\Phi$ is the form map and $\mathcal F$ its codomain. The metrics $d_{\bar Q}$ and $d_{\mathcal F}$ measure state and form distance respectively.

The symbols $\mathscr C$ and $\mathscr V$ are reserved for coherence and relational variation. $\mathcal M$ denotes a constraint manifold and $\mathcal A$ an attractor. $\mathrm{Poss}$ denotes an admissible possibility space and $\mathcal K_{\mathrm{act}}$ an actualization kernel. $U_g$ denotes a representation of a group element and $P_G$ the Reynolds projector. $\mathscr C_G$ denotes finite-group symmetry coherence. The symbol $K$ is reserved for complex conjugation in the wave chapters and $D$ for Platonic duality. $V_\ell^{\mathbb R}$ denotes the real spherical-harmonic representation and $\mathcal H_{\mathbb C}$ the complex wave Hilbert space.

In the ontology chapters, $P_{\mathrm{phys}}$ denotes a candidate physical state description, $E_{\mathrm{exp}}$ the experiential state posited by an ontology, and $Z_\varphi$ the empirically inferred phenomenological-evidence latent. These three symbols are not interchangeable. $\Pi_{\mathrm{ext}}$ and $\Pi_{\mathrm{int}}$ denote exterior and interior projections in a dual-aspect model. The model families $M_P$, $M_Z$, and $M_R$ denote physical-only, phenomenological-evidence, and constrained relational models respectively.

\section*{Claim ledger}

The exact mathematical results include the projector identities $P_G^2=P_G$ and $P_G^\dagger=P_G$; the orthogonal coherence decomposition $\mathscr C_G+\mathscr E_G=1$; character multiplicity as the dimension of the finite-group fixed space; the parity statements separating $SO(3)$ and $O(3)$ invariants; the fact that geometry does not determine a unique dynamics; commutation of scalar conjugation with rotations and group projection; preservation of density and symmetry coherence under conjugation together with reversal of the standard phase current; the general non-equivalence of conjugation and Platonic duality; the exact quartic channel relations in $V_5$; the reduction $F_4=\alpha K_0+\beta K_2$; the quartic minimum set $\mathcal M_4=\{B_2=0\}\cap S^{10}$; the rank-five certificate at $A_\star$; the local five-dimensional minimum manifold; its decomposition into three rotational and two shape directions; the dimension count showing four primitive sixth-order invariant directions beyond $I_2\operatorname{Inv}_4$; and the positive-definite shape Hessian of the explicit sextic invariant $J_{4,8}$ at $A_\star$.

The model-specific results include the observed octahedral $\ell=4$ and icosahedral $\ell=6$ asymptotic branches in the stated nonlinear spherical PDE and parameter regime, the associated approximate transverse Hessian gaps, and the observed sensitivity of relaxation time to quadratic coupling. These are not promoted to universal theorems.

The sixth-order $V_5$ result is an algebraic capability result. It proves that the invariant algebra contains structure able to resolve the quartic shape-flat plane. It does not yet prove the coefficient with which that structure appears in the complete sixth-order center-manifold reduction of the chosen PDE.

The A0--A7 ladder, nested physical sufficiency test, partial information decomposition protocol, complexity-controlled prequential comparison, and defeater rules are methodological proposals. They are not empirical results until applied to data.

Relational physicalism, relational idealism, and relational dual-aspect realism are ontological hypotheses. None is presented as a theorem forced by the mathematical work.
